% ============================== CI-4 ==============================
\reading{CI-4}{Global Financial Stability Report, Chapter 2: Geopolitical Risk}{STRIKE FIRST}{2}

\begin{cikeybox}[title={The six objectives, in the spine's own words}]
\begin{enumerate}[label=\textbf{\alph*.}]
\item Identify and explain the key channels through which geopolitical risk influences asset
prices and financial stability and discuss policy measures to address potential consequences.
\item Analyze how global geopolitical risk events affect countries, sectors, and asset classes
differently, including their varying impacts on commodity importers and exporters.
\item Assess the stock market responses to different domestic and global geopolitical risk
events, including Russia's invasion of Ukraine and the U.S.--China trade tensions, and explain
the cross-border effects of these events.
\item Evaluate the effects of geopolitical shocks on sovereign risk premiums.
\item Explain how investors price geopolitical risk in equity and option markets, including
changes in downside risk and tail-risk premiums.
\item Describe the implications of geopolitical risk exposure for the resilience of banks and
nonbank financial institutions, including effects on equity valuations, lending, and
investment fund flows.
\end{enumerate}
\greyline{Source: Global Financial Stability Report, Chapter 2, IMF (April 2025). The class
notes label this Reading 103; it is CI-4 on the 2026 spine. Six objectives makes this the
longest reading in the area, and it is also the most number-heavy.}
\end{cikeybox}

% ------------------------------------------------------------------
\lo{CI-4 a}{Identify and explain the key channels through which geopolitical risk influences asset prices and financial stability and discuss policy measures to address potential consequences.}

\begin{cidefbox}[title={Geopolitical risk}]
Disruptions from \term{wars, terrorism, trade conflicts (tariffs/sanctions)} that impact
financial markets, trade, and growth. Even the \term{threat} of tensions can reduce asset
prices and shift capital flows.
\end{cidefbox}

Geopolitical risk affects asset prices through two primary channels.

\begin{cifigblock}
{\centering
\begin{tikzpicture}[
  font=\sffamily\scriptsize,
  ch/.style={draw=FRMRule, fill=PaleNavy, text width=4.5cm, align=left,
             minimum height=1.25cm, inner sep=5pt, rounded corners=1pt},
  eff/.style={draw=FRMRule, fill=PaleGold, text width=4.6cm, align=left,
              inner sep=5pt, rounded corners=1pt},
  arr/.style={-{Stealth[length=4.5pt]}, FRMGold, line width=0.8pt}
]
\node[ch] (c1) at (0,0) {\textbf{Economic channel}\\[2pt] disrupts trade, capital flows and
  production capacity};
\node[eff, right=1.5cm of c1] (o1) {\textbf{Leads to}
  \begin{itemize}[leftmargin=10pt, topsep=2pt, itemsep=1pt]
  \item fall in equities and bonds
  \item rise in commodity prices (oil, grains)
  \item higher sovereign debt (military spending, lower revenues)
  \end{itemize}};
\node[ch, below=1.6cm of c1] (c2) {\textbf{Market sentiment channel}\\[2pt] driven by
  uncertainty, fear and risk aversion};
\node[eff, right=1.5cm of c2] (o2) {\textbf{Even without actual conflict, leads to}
  \begin{itemize}[leftmargin=10pt, topsep=2pt, itemsep=1pt]
  \item lower valuations
  \item higher volatility
  \end{itemize}};
\draw[arr] (c1) -- (o1) node[midway, above, font=\sffamily\scriptsize\color{FRMGold},
  fill=white, inner sep=1.5pt] {produces};
\draw[arr] (c2) -- (o2) node[midway, above, font=\sffamily\scriptsize\color{FRMGold},
  fill=white, inner sep=1.5pt] {produces};
\end{tikzpicture}\par}
\figcap{Figure CI-4.1. Cause on the left, effect on the right. The discriminating feature is
the trigger, not the outcome: the economic channel needs the disruption to actually happen,
while the market sentiment channel fires on the \emph{threat} alone. That is why the reading
insists that even the threat of tensions moves prices, and it is why the second channel is
the one that operates in markets far from the conflict.}
\end{cifigblock}

\subsection*{Policy measures}

\begin{enumerate}[label=\textbf{\alph*)}]
\item \term{At country level}: strengthen regulation, maintain \term{fiscal and monetary
flexibility}.
\item \term{Financial institutions}: maintain \term{capital and liquidity buffers}.
\item \term{Goal}: to improve resilience to shocks.
\end{enumerate}

% ------------------------------------------------------------------
\lo{CI-4 b}{Analyze how global geopolitical risk events affect countries, sectors, and asset classes differently, including their varying impacts on commodity importers and exporters.}

\subsection*{Impact of geopolitical risk across countries, sectors, and asset classes}

\begin{center}
\renewcommand{\arraystretch}{1.25}
\begin{tabularx}{\textwidth}{@{}p{2.4cm} >{\raggedright\arraybackslash}X
                              >{\raggedright\arraybackslash}X@{}}
\toprule
\rowcolor{FRMSoft}
{\sffamily\bfseries\color{FRMNavy}Cut} &
{\sffamily\bfseries\color{FRMGreen}Gainers / appreciate} &
{\sffamily\bfseries\color{FRMRed}Losers / decline}\\
\midrule
\term{Countries}
& Commodity exporters benefit from higher prices
& Conflict countries take the most severe damage: physical destruction, sanctions, trade
disruptions. Commodity importers suffer currency depreciation and higher costs\\
\addlinespace
\term{Sectors}
& Defence and energy, from increased military spending and higher commodity prices
& Banking, retail and utilities, from weakened demand and tighter credit\\
\addlinespace
\term{Asset classes}
& Commodities: oil and gold usually appreciate on supply fears and safe-haven demand.
Safe-haven sovereign bonds appreciate, so their yields fall
& Equities generally decline: a standard event causes about a 3\% drop, and a severe shock can
trigger declines of up to 9\%\\
\bottomrule
\end{tabularx}
\end{center}
\figcap{Table CI-4.1. The impact is uneven, and it depends on economic structure and exposure.
Reading across the row rather than down the column is what the objective asks for: the same
event that lifts a commodity exporter's terms of trade is the one that depreciates a commodity
importer's currency.}

\subsection*{Severity of impact}

Severity depends on the type and duration of the event.

\begin{itemize}
\item \term{Wars and military conflicts}: largest and long-lasting impact.
\item \term{Sanctions and trade disputes}: relatively smaller impact.
\end{itemize}

% ------------------------------------------------------------------
\lo{CI-4 c}{Assess the stock market responses to different domestic and global geopolitical risk events, including Russia's invasion of Ukraine and the U.S.--China trade tensions, and explain the cross-border effects of these events.}

\subsection*{General stock market responses}

Stock markets usually react mildly, but large or global shocks amplify the impact
significantly. The effect depends on uncertainty plus investor risk aversion.

\begin{itemize}
\item Average impact: approximately 0.3\% fall in returns after a country-specific shock.
\item Large regional shocks (above 2 SD): up to \term{7 times} larger losses.
\item Global shocks: up to \term{20 times} larger impact, with longer persistence.
\end{itemize}

\begin{cifigblock}
{\centering
\begin{tikzpicture}
\begin{axis}[
  width=11.5cm, height=4.3cm,
  xbar, bar width=13pt,
  xmode=log, log basis x=10,
  xmin=0.8, xmax=60,
  xtick={1,2,5,10,20,50},
  xticklabels={1$\times$,2$\times$,5$\times$,10$\times$,20$\times$,50$\times$},
  xlabel={\sffamily\scriptsize Size of the equity market reaction, as a multiple of the
          country-specific baseline (log scale)},
  symbolic y coords={Global shock, Large regional shock, Country-specific shock},
  ytick=data,
  tick label style={font=\sffamily\scriptsize},
  label style={font=\sffamily\scriptsize},
  nodes near coords,
  nodes near coords style={font=\sffamily\scriptsize\color{FRMNavy},
     fill=white, inner sep=1.5pt, anchor=west, xshift=3pt},
  point meta=explicit symbolic,
  xmajorgrids, grid style={FRMRule, dashed},
  axis line style={FRMRule},
  enlarge y limits=0.30
]
\addplot[fill=FRMNavy!70, draw=FRMNavy] coordinates {
  (1,{Country-specific shock})    [baseline: about 0.3\% fall]
  (7,{Large regional shock})      [up to 7$\times$, above 2 SD]
  (20,{Global shock})             [up to 20$\times$, longer persistence]
};
\end{axis}
\end{tikzpicture}\par}
\figcap{Figure CI-4.2. The axis is logarithmic on purpose. On a linear scale the baseline
reaction is invisible next to the global one, and the point of the objective is precisely that
the three are orders of magnitude apart rather than merely different. Note that the
multipliers are ceilings, ``up to'', not averages, and that only the global case is described
as persistent.}
\end{cifigblock}

\subsection*{Volatility}

The CBOE Volatility Index (VIX) spikes after shocks, which leads to higher uncertainty, and
tail risk increases (higher probability of extreme losses).

Military conflicts have the strongest impact on equity markets, especially in emerging
economies. Effects spread globally through trade and financial linkages, with nearby or
economically connected countries affected most.

\subsection*{The two cases}

\begin{center}
\renewcommand{\arraystretch}{1.25}
\begin{tabularx}{\textwidth}{@{}p{2.9cm} >{\raggedright\arraybackslash}X
                              >{\raggedright\arraybackslash}X@{}}
\toprule
\rowcolor{FRMSoft}
&{\sffamily\bfseries\color{FRMNavy}Russia--Ukraine war (2022)} &
{\sffamily\bfseries\color{FRMNavy}US--China trade tensions (2018--2024)}\\
\midrule
\term{Type of event} & Military conflict & Tariff measures\\
\addlinespace
\term{Where it landed hardest} & Broad market declines, with Europe more affected due to
strong economic ties and energy dependence &
China, where markets fell by about 4\% on average and nearly 8\% after the May 2019 US
tariffs\\
\addlinespace
\term{Direct exposure effect} & Firms with direct exposure saw valuation drops of about 2.5\%
within a week &
US markets less affected: average declines of around 1.3\%, and about 1.6\% to 1.8\% following
China's retaliatory tariffs\\
\addlinespace
\term{Sectoral split} & Defence and energy firms \emph{with exposure} declined initially,
while similar firms in \emph{other countries} gained on expectations of higher spending and
energy prices & \\
\addlinespace
\term{Since} & Overall exposure to Russia has reduced since 2022 & \\
\bottomrule
\end{tabularx}
\end{center}
\figcap{Table CI-4.2. The two cases are the reading's own contrast pair, and the row that
carries the most weight is the sectoral split. Being in the defence sector was not enough to
gain from the invasion; being in the defence sector \emph{outside} the exposed economies was.
Sector and geography have to be read together.}

% ------------------------------------------------------------------
\lo{CI-4 d}{Evaluate the effects of geopolitical shocks on sovereign risk premiums.}

\begin{itemize}
\item Sovereign risk premiums tend to rise following major geopolitical shocks, particularly
in emerging markets and commodity-importing countries.
\item Sovereign risk premiums are commonly assessed using \term{sovereign credit default swap
(CDS) spreads}. These spreads tend to rise sharply after military conflicts because investors
become more concerned about fiscal stability and debt sustainability.
\end{itemize}

Emerging markets are more vulnerable due to weaker fundamentals.

\begin{cikeybox}[title={The yield split, which runs in both directions at once}]
\begin{itemize}
\item \term{Vulnerable markets}: yields \textcolor{FRMRed}{\bfseries rise} as investors demand
higher risk premiums.
\item \term{Safe havens}: yields \textcolor{FRMGreen}{\bfseries fall} (e.g., U.S., Japan,
Germany) due to a ``flight to quality''.
\end{itemize}
\vspace{3pt}
Major risk drivers: a) high debt-to-GDP, b) low foreign reserves, c) weak institutional
quality.
\end{cikeybox}

% ------------------------------------------------------------------
\lo{CI-4 e}{Explain how investors price geopolitical risk in equity and option markets, including changes in downside risk and tail-risk premiums.}

Investors price geopolitical risk into both equity and option markets through risk premiums.
Assets that \emph{lose} value when geopolitical risks rise require \term{positive} risk
premiums, while assets that tend to \emph{gain} value during periods of elevated geopolitical
stress often exhibit \term{negative} premiums.

\subsection*{Equity markets and GPR betas}

\term{GPR beta} measures a stock's sensitivity to geopolitical shocks.

\begin{cifigblock}
{\centering
\begin{tikzpicture}[
  font=\sffamily\scriptsize,
  bx/.style={draw=FRMRule, text width=4.3cm, align=left, inner sep=5pt,
             minimum height=1.25cm, rounded corners=1pt}
]
\draw[FRMRule, line width=1.0pt] (-0.2,0) -- (5.35,0);
\node[font=\sffamily\bfseries\scriptsize\color{FRMNavy}, fill=white, inner sep=2pt]
  at (2.6,0) {GPR beta: the sign splits the sectors};

\node[bx, fill=PaleGreen] (pos) at (2.6,1.25) {\textbf{Positive beta: gainers}\\[2pt]
  energy and defence, which benefit from rising commodity prices and increased military
  spending};
\node[bx, fill=PaleRed] (neg) at (2.6,-1.25) {\textbf{Negative beta: losers}\\[2pt]
  consumer-oriented sectors, which suffer as uncertainty weakens demand};

\node[bx, fill=PaleGold, text width=5.0cm] (opt) at (8.9,1.25)
  {\textbf{Option market pricing}\\[2pt] geopolitical risk is priced through put option
  premiums, which are downside protection};
\node[bx, fill=PaleNavy, text width=5.0cm] (tail) at (8.9,-1.25)
  {\textbf{Tail risk}\\[2pt] put premiums spike during shocks, and demand for ``deep
  out-of-the-money'' puts rises most sharply};

\node[align=center, font=\sffamily\scriptsize\itshape\color{FRMGrey}, fill=white,
      inner sep=2pt] at (2.6,-2.55)
  {after the 2022 invasion of Ukraine the energy sector's premium\\
   shifted from \textbf{negative to positive} as it became a critical hedge};
\end{tikzpicture}\par}
\figcap{Figure CI-4.3. Equity pricing on the left, option pricing on the right. The sign
convention is the thing to hold: a \emph{positive} risk premium is demanded on assets that
\emph{lose} value in a shock, so a hedge asset carries a \emph{negative} premium. That is what
makes the energy note at the bottom a genuine reversal rather than a detail, and it is the one
dated sign flip the reading gives you.}
\end{cifigblock}

% ------------------------------------------------------------------
\lo{CI-4 f}{Describe the implications of geopolitical risk exposure for the resilience of banks and nonbank financial institutions, including effects on equity valuations, lending, and investment fund flows.}

\subsection*{Effects on banks and nonbank financial institutions}

\begin{enumerate}[label=\textbf{\alph*)}]
\item Geopolitical shocks increase market, credit, and liquidity risks due to volatility and
uncertainty.
\item Selloffs reduce asset and collateral values, leading to tighter lending conditions.
\item Capital outflows and sanctions amplify stress; cyber risks also increase.
\item As of 2024, about 8\% of cross-border bank claims and 13\% of equity fund assets are
tied to high-risk geopolitical markets.
\item Banks and funds have reduced direct exposure to Russia and Ukraine after the 2014 and
2022 events.
\end{enumerate}

\subsection*{Effects on equity valuations and lending}

Geopolitical shocks weaken borrower creditworthiness, leading to:

\begin{enumerate}[label=\textbf{\alph*)}]
\item sharp declines in cross-border credit lending,
\item lower asset valuations and bank equity levels,
\item difficulty in refinancing foreign debt during heightened uncertainty (rollover risk).
\end{enumerate}

\subsection*{Effects on investment fund flows}

Funds with higher exposure to conflict regions face redemptions, lower returns, and valuation
pressure, with stronger impacts observed during major conflicts like the Russia--Ukraine war
compared to more localised events like US--China trade tensions.

\begin{cifigblock}
{\centering
\begin{tikzpicture}[
  font=\sffamily\scriptsize,
  st/.style={draw=FRMRule, fill=PaleRed, text width=2.55cm, align=center,
             minimum height=1.15cm, inner sep=3pt, rounded corners=1pt},
  arr/.style={-{Stealth[length=4.5pt]}, FRMRed, line width=0.8pt}
]
\node[st, fill=FRMSoft] (s0) at (0,0) {geopolitical shock};
\node[st, right=6mm of s0] (s1) {market, credit and liquidity risks rise};
\node[st, right=6mm of s1] (s2) {selloffs cut asset and collateral values};
\node[st, right=6mm of s2] (s3) {tighter lending conditions, cross-border lending falls};
\node[st, right=6mm of s3, fill=PaleGold] (s4) {redemptions, lower returns, valuation
  pressure on funds};
\draw[arr] (s0) -- (s1);
\draw[arr] (s1) -- (s2);
\draw[arr] (s2) -- (s3);
\draw[arr] (s3) -- (s4);
\node[below=4mm of s2, font=\sffamily\scriptsize\itshape\color{FRMGrey}, fill=white,
      inner sep=2pt, text width=8.5cm, align=center]
  {rollover risk sits here: foreign debt becomes hard to refinance exactly when
   collateral values and bank equity are falling};
\end{tikzpicture}\par}
\figcap{Figure CI-4.4. The transmission runs from price to collateral to credit to flows, and
it is a chain rather than four separate effects. The last link is where banks and nonbanks
part company: banks feel it as tighter lending conditions, funds feel it as redemptions, and
the reading says the fund channel bites hardest in major conflicts rather than localised trade
disputes.}
\end{cifigblock}

\begin{cikeybox}[title={Closing position of the reading}]
Overall, geopolitical shocks can threaten financial stability, especially if institutions lack
adequate capital and liquidity buffers. Strengthening resilience requires better risk
management, diversification, and strong regulatory frameworks.
\end{cikeybox}
